\documentclass[11pt,a4paper]{article}
\usepackage[T1]{fontenc}
\usepackage[utf8]{inputenc}
\usepackage{lmodern,microtype}
\usepackage[margin=25mm,top=23mm,bottom=24mm,headheight=14pt]{geometry}
\usepackage{amsmath,amssymb,amsthm,booktabs,array,xcolor,fancyhdr}
\definecolor{ink}{HTML}{143649}
\definecolor{accent}{HTML}{0B625E}
\definecolor{quiet}{HTML}{4C626A}
\usepackage[colorlinks=true,linkcolor=accent,urlcolor=accent,citecolor=accent,hypertexnames=false,pdftitle={Deleting mass on the integers},pdfauthor={Shaitetsu}]{hyperref}
\newtheorem{theorem}{Theorem}[section]
\newtheorem{proposition}[theorem]{Proposition}
\newtheorem{corollary}[theorem]{Corollary}
\newtheorem{lemma}[theorem]{Lemma}
\newtheorem*{remark}{Remark}
\newcommand{\E}{\mathbb{E}}
\newcommand{\N}{\mathbb{N}}
\newcommand{\R}{\mathbb{R}}
\newcommand{\supp}{\operatorname{supp}}
\pagestyle{fancy}\fancyhf{}\fancyhead[L]{\small\color{quiet}Deleting mass on the integers}
\fancyhead[R]{\small\color{quiet}Arithmetic and combinatorial applications}\fancyfoot[C]{\thepage}
\renewcommand{\headrulewidth}{0.3pt}\setlength{\emergencystretch}{2em}\raggedbottom
\begin{document}
\thispagestyle{plain}
\noindent{\small\sffamily\bfseries\color{accent} NUMBER THEORY\quad /\quad EXTREMAL COMBINATORICS}\par
\vspace{9pt}\noindent{\LARGE\bfseries\color{ink} Deleting mass on the integers}\par
\vspace{4pt}\noindent{\large GCD kernels, prime-factor distance and $p$-adic selection}\par
\vspace{9pt}\noindent{\small Shaitetsu\quad$\cdot$\quad 21 November 2026}\par
\vspace{9pt}\hrule\vspace{10pt}
\noindent\textbf{Abstract.}
A single directed network represents every maximum-perimeter subset of a finite weighted ultrametric space, including ties at every cardinality. This gives an affirmative answer to a published question of Grinberg and Petrov arising from Bhargava's $p$-orderings. For two normalized greatest-common-divisor kernels, stationary probabilities give opposite support decisions on all primes in $(N/6,N/5]$ for every sufficiently large $N$. We also determine the leading coefficient for the original distance energy on the graph of integers linked by prime multiplication. Applications to equal pairwise least common multiples give an exact extremal family on semiprimes. For equal pairwise greatest common divisors we express an arithmetic growth rate in terms of sunflower capacity. Each transfer identifies the arithmetic quantity preserved by its graph or set representation.
\section{The part of the graph argument that transfers}
In the graph manuscript motivating this paper, a probability vector maximizes $p^{\mathsf T}Kp$, with $K_{ii}=2$ and $K_{ij}$ the distance between vertices for $i\ne j$. Stationarity and deletion force a minimum amount of mass at certain vertices of a hypothetical counterexample. A comparable argument can be used on integers after specifying the arithmetic kernel, what deleting an integer means, and how the kernel on the remaining indices changes.
\begin{lemma}[Exact removal cost]\label{lem:remove}
Let $K$ be a real symmetric matrix on a finite index set. Let $p$ be a maximizer or a minimizer of $p^{\mathsf T}Kp$ on the probability simplex, and set $Q=p^{\mathsf T}Kp$. At each index of positive mass, $(Kp)_i=Q$. If $p_i=t<1$, let $q$ be the probability on the surviving indices obtained by renormalizing after deleting $i$. With the original kernel restricted to them,
\begin{equation}\label{eq:remove}
 q^{\mathsf T}Kq-Q=(K_{ii}-Q)\frac{t^2}{(1-t)^2}.
\end{equation}
\end{lemma}
\begin{proof}
Transfers between any two positive coordinates are allowed in both directions. The first derivative is twice the difference of their potentials, so those potentials agree. Their $p$-weighted average is $Q$. Expanding $(p-t e_i)^{\mathsf T}K(p-t e_i)/(1-t)^2$ proves the identity for $t>0$. At $t=0$ both sides vanish.
\end{proof}
For a fixed arithmetic kernel, deletion means passing to a principal submatrix. At a zero coordinate of a minimizer, the potential is at least $Q$, by the one-sided derivative. If the minimizer is unique, the optimized value stays equal after a deletion exactly when the deleted coordinate has zero mass. The formula above gives the cost of the particular renormalized competitor; the surviving vertices may be reoptimized. In the final combinatorial applications the variables are counts or indicators. There the relevant operation is a one-unit exchange, and its gain must be derived separately.

\section{A linkage representation for Bhargava's greedoid}
\label{sec:bhargava-linking}
Grinberg and Petrov asked whether the Bhargava greedoid of every finite ultra triple is a linking greedoid \cite[Question~29]{gp}. The question appears as Question~6.5 in the June 2026 revision \cite{gprev}. The Gaussian-elimination representation over sufficiently large fields was proved in \cite{grin}. Here we construct a directed network simultaneously representing every cardinality. We use the directed, vertex-disjoint linkage definition in \cite[Sec.~IV.2.15, p.~53 and Sec.~IX.4, p.~137]{kls}.

\begin{theorem}[An affirmative answer to Question 29]\label{thm:gp}
Let $E$ be a finite set, let $w:E\to\mathbb R$, and let $d$ be a real symmetric
function on pairs of distinct elements satisfying the ultrametric inequality
$d(x,z)\le\max\{d(x,y),d(y,z)\}$ for distinct $x,y,z$. Define
\[
 P(S)=\sum_{x\in S}w(x)+\sum_{\{x,y\}\subseteq S}d(x,y).
\]
There is a finite directed acyclic graph with distinct ordered sources
$s_1,\ldots,s_{|E|}$ and separate terminals indexed by $E$ such that
$S\subseteq E$, $|S|=k$, attains the maximum of $P$ among the $k$-sets
if and only if $k$ pairwise vertex-disjoint directed paths join
$\{s_1,\ldots,s_k\}$ to $S$. Their pairing is unrestricted.
Reversing all arcs gives paths from the chosen elements of $E$ to the ordered target prefix, as in the definition of a linking greedoid.
\end{theorem}

\paragraph{Local increments.}
For $|E|\ge 2$ build a tree whose root is $E$.
At an internal cluster $v$, let
$h_v=\max_{x\ne y\in v}d(x,y)$ and take as children the classes of
$x=y$ or $d(x,y)<h_v$. The ultrametric inequality gives transitivity,
and a pair attaining $h_v$ lies in two separate classes.
All distances between different children equal $h_v$.
Each internal child $u$ satisfies $h_u<h_v$.
For an internal $v$ define
\[
 A_v(k)=\max_{\substack{S\subseteq v\\|S|=k}}
             \biggl(P(S)-h_v\binom{k}{2}\biggr),
 \qquad
 \alpha_v(i)=A_v(i)-A_v(i-1).
\]
At the leaf $\{x\}$ put $A(0)=0$, $A(1)=\alpha(1)=w(x)$.
For each child $u$ of $v$ form the list of labels
\[
 b_{vu}(i)=\alpha_u(i)-(h_v-h_u)(i-1)
 \quad(1\le i\le |u|,\ u\text{ internal}),\qquad
 b_{v\{x\}}(1)=w(x).
\]
We prove from the leaves up that the nonincreasing rearrangement of all the
child labels at $v$ is precisely $\alpha_v(1),\ldots,\alpha_v(|v|)$.
Indeed, fixing the child counts $k_u$ makes the best normalized value
\[
 \sum_{u}\sum_{i=1}^{k_u} b_{vu}(i).
\]
Cross-child pairs disappear in the normalization and each internal
child contributes $A_u(k_u)-(h_v-h_u)\binom{k_u}{2}$.
By induction $\alpha_u$ is nonincreasing. Since the height gap is positive,
each individual child list is strictly decreasing. Therefore the largest
sum for $\sum k_u=k$ is the sum of the $k$ largest labels in the union.
Every choice of these labels, including any choice at the tie for the last
place, consists of an initial segment of each child list.
This proves the rearrangement claim. It also gives its equality condition:
a specified set $S$ is optimal inside $v$ precisely when each
$S\cap u$ is optimal at its own cardinality in $u$ and its child prefixes
select a multiset of the $k$ largest labels.

\paragraph{The network and its proof.}
For each cluster $v$ create vertices $(v,i)$, one for each $1\le i\le |v|$.
For every child $u$ put in all arcs
\[
        (v,j)\longrightarrow(u,i)
        \qquad\text{such that}\qquad
        \alpha_v(j)=b_{vu}(i).
\]
Put an arc from every singleton vertex to its own ground-set terminal.
Use the root vertices in increasing index as sources.
For every $v$ and $k$ we prove that the terminal sets reachable by
vertex-disjoint paths from the first $k$ vertices of $v$ are exactly the
optimal $k$-sets inside $v$. The assertion is immediate at a singleton
and at $k=0$. At an internal cluster the first steps of disjoint paths
occupy $k$ distinct child vertices bearing exactly the $k$ largest labels.
Strict decrease within each child forces its occupied indices to be a
prefix. The remaining paths in each child certify local optimality by
induction. The equality condition gives optimality in $v$. Conversely,
for an optimal set the selected child prefixes carry exactly the top $k$
labels. For each label choose a bijection from the corresponding occupied
parent positions to the selected child positions bearing that label.
All required arcs exist. Complete the paths using the inductive linkages
in the pairwise disjoint child subtrees. The result holds at the root,
where subtracting $h_E\binom k2$ changes no fixed-cardinality optimum.
The singleton construction works for $|E|=1$, and the empty graph suffices
for $E=\varnothing$.

\paragraph{Size and a matroid consequence.}
Writing $n=|E|$, the graph has $O(n^2)$ vertices and arcs.
The tree has at most $2n-1$ clusters. At an internal $v$, the multiplicity
of each label is bounded by the number $t_v$ of children because a child
list is strictly decreasing. The number of arcs at $v$ is the sum of the
squares of the multiplicities, at most $|v|t_v$. Now
$\sum_v t_v\le 2n-2$ for $n\ge1$. For each $k$ the maximizing $k$-sets
are the bases of a gammoid: use the same graph with the first $k$ sources
and allow subsets of those sources for independent sets. Its rank is $k$
because an optimal $k$-set exists.

\paragraph{A $p$-adic example.}
For distinct integers take $d(x,y)=-v_p(x-y)$, where $v_p$ is the exponent
of $p$ in a nonzero integer, and take $w=0$. Each selected integer contributes one unit of mass. If $m_{j,a}$ integers in $S$ fall in the residue class $a$ modulo $p^j$, then
\[
  \sum_{\{x,y\}\subseteq S}v_p(x-y)
  =\sum_{j\ge1}\sum_{a\bmod p^j}\binom{m_{j,a}}{2}.
\]
Only finitely many summands are nonzero. Removing one chosen integer from a class of size $m$ removes $m-1$ collisions at that level. Adding it to a class of size $b$ creates $b$ collisions. The labels in the proof organize these one-unit gains across all nested classes. Here $P(S)$ is the negative of
the $p$-adic valuation of the Vandermonde product of $S$. For $p=2$ and
$E=\{0,1,2,3\}$ the two parity classes each supply parent labels $(0,-1)$.
The first two sources can reach precisely pairs containing one even and one
odd element. Thus a single directed graph records all minimizing sets,
including ties, for all product sizes. The network proof treats one
ultrametric at a time.

In the four-point example a one-step bipartite matching would fail. Since every singleton is optimal, its first source must see all four points. Each forbidden same-parity pair then forces the second source to miss both its points. The second source misses all four, contradicting the feasibility of every mixed pair. Intermediate vertices in the linkage construction carry necessary information.

\paragraph{Exact finite checks.}
For each test we enumerated every subset, computed its perimeter, and compared the fixed-size optimality decision with a unit-capacity maximum flow in the constructed graph, splitting each vertex to enforce vertex disjointness. The tests included all 60 ultrametrics on four labelled elements taking off-diagonal values in $\{-1,0,1\}$ with each of the 16 binary weight assignments, 80 generated examples of sizes six to nine (integer distances and weights, seed 290341), and the eight-element set $\{0,1,2,4,6,9,10,13\}$ with $d=-v_3$ and zero weights. All 34,816 subset decisions in these 1,041 instances agreed. There were 2,232 instance/cardinality pairs with tied maximizers. A separate test of all 256 subsets in Example~27 of \cite{gp} also agreed. These computations check the construction independently of the inductive proof.

\section[Removing the r nearest predecessors]{Removing the \(r\) nearest predecessors}
\label{sec:removed}
Krachun and Mirgalimova proved that the maximizing sets for the
\(r\)-removed perimeter form a strong greedoid
\cite[Definitions~2--3, 11 and Theorem~19]{removed}.
Their perimeter adds the weight of each newly selected point \(x\) and
the sum of its distances to its predecessors after discarding the \(r\)
smallest distances. If there are at most \(r\) predecessors this
distance contribution is zero. Ties give the same numerical sum.
We show that the maximizing sets also admit a single directed
linkage representation for all cardinalities, for each fixed \(r\ge0\).

Let \(v_0\) be the root of the cluster tree, write \(k_u=|S\cap u|\),
and let \(\delta_u=h_{\mathrm{parent}(u)}-h_u>0\) for nonroot internal
clusters \(u\). Put \(g_r(k)=\binom{\max(k-r,0)}{2}\), so
\(g_r(k+1)-g_r(k)=\max(k-r,0)\).
The following identity is independent of the ordering of \(S\):
\[
 P_r(S)=\sum_{x\in S}w(x)+h_{v_0}g_r(|S|)
           -\sum_{\substack{u\ne v_0\\u\ {\rm internal}}}
                           \delta_u\,g_r(k_u).
\]
Indeed, for a new point \(x\) with \(m\) predecessors, its contribution
from the root height is \(h_{v_0}(m-r)_+\), where \(a_+=\max(a,0)\).
For any proper internal cluster \(u\ni x\), the \(m_u\) predecessors in
\(u\) are all strictly nearer to \(x\) than any predecessor outside \(u\).
Hence exactly \((m_u-r)_+\) of them survive, producing the correction
\(-\delta_u(m_u-r)_+\). Every surviving distance is its root height minus
the height gaps on the common chain, so these contributions are exact.
Summing over arrivals gives the displayed identity. Singletons and the
empty set need only their weight sum.

At each internal \(v\), now subtract \(h_v g_r(k)\) from its local
\(r\)-removed perimeter when computing the optimum \(A_v(k)\); define
\(\alpha_v(i)=A_v(i)-A_v(i-1)\) as before. The parent label of an
internal child \(u\) becomes
\[
 b_{vu}(i)=\alpha_u(i)-\delta_u(i-1-r)_+.
\]
Singletons retain their weight. By induction each child list is weakly
decreasing. Its first \(k_u\) entries give the optimal contribution from
that child, so an optimal allocation is again a selection of the \(k\)
largest labels. Such an allocation can always be taken in prefixes by
ordering tied labels within a child in increasing index. Conversely the
chosen prefixes of any optimal count vector must have the top label
multiset. Thus the parent increments are still the sorted union of the
child labels. This proves the induction on weak decrease.

Use the same graph, with arcs wherever parent increments and child labels
coincide. The one change in the path proof concerns the indices occupied
in a child \(u\). All indices with \(b_{vu}\) strictly above the selection
threshold are occupied. Within a tie, \(b_{vu}(i)=b_{vu}(j)\) for \(i<j\)
implies \(\alpha_u(i)=\alpha_u(j)\), because both terms in the formula
for \(b_{vu}\) are weakly decreasing. Vertices of \(u\) with the same
\(\alpha_u\) have identical outgoing neighbors. Permuting those tied
indices therefore changes neither the available incoming arcs from the
used parent positions nor the reachable terminals. We may move the
occupied indices to the earliest available positions in their tie block,
so that \(u\) receives a prefix. Applying the induction in each child
shows every linkage yields an optimal set. Conversely an optimal set's
child prefixes carry the top multiset, hence may be matched label by
label and extended by the local linkages. This proves the representation
also for \(r>0\), including the rankwise gammoid conclusion.
A given label occurs at most \(r+1\) times in a child list because its
penalty strictly increases after position \(r+1\). Consequently the
same arc count gives \(O((r+1)n^2)\) arcs and \(O(n^2)\) vertices.

For example take \(E=\{0,2,4,6,1\}\) with \(d(x,y)=-v_2(x-y)\),
\(w=0\), \(r=1\). Each pair has perimeter zero. For a triple only the last
point can contribute a distance, and it keeps the larger of its two
distances. This equals the diameter of the triple by ultrametricity.
Thus a triple containing \(1\) is optimal with value zero, while every
all-even triple has negative value.

A separate enumeration compared scores computed by removing the $r$ nearest predecessors in an explicit ordering with the vertex-split flow test. The scores were also checked with the arrival order reversed. All 61,440 subset comparisons agreed across 660 (space, $r$) pairs on at most eight elements.

\section{The optimizers of two small Gál sums}\label{sec:gal}

De la Bretèche, Munsch and Tenenbaum study the kernels
\[
 K^{T}(a,b)=\frac{\gcd(a,b)}{\sqrt{ab}},\qquad
 K^{V}(a,b)=\frac{\gcd(a,b)}{a+b}.
\]
For either kernel write $K_N=(K(a,b))_{1\le a,b\le N}$ and set
\[
 \lambda_N=\min_{w\ge0,\ \sum w_i=1}w^{\mathsf T}K_Nw,\qquad
 C_N=1/\lambda_N,\qquad C_0=0 .
\]
These symbols refer to the chosen kernel. The published quantities $T_N,V_N$ equal $N\lambda_N$ for $K^T,K^V$, respectively. Theorem~1.1 of \cite{dmtgal} gives, with $\eta\simeq0.16656$,
\[
 (\log N)^\eta\ll V_N\le T_N/2\ll
 (\log N)^\eta(\log\log N)^3.
\]
We use its lower bound. We determine which integers receive mass in several prime intervals in the unrestricted optimizers and how the optimum value changes when a single integer is added.

\begin{theorem}[Prime intervals in the global optimizers]\label{thm:gal-primes}
For all sufficiently large $N$ the following assertions hold simultaneously for all primes in the indicated ranges.
\begin{enumerate}
\item In the $K^T$ minimizer, every prime $N/5<p\le N$ has positive mass. Every prime $N/6<p\le N/5$ has zero mass. Every prime $N/7<p\le N/6$ again has positive mass, and $2p$ has zero mass.
\item In the $K^V$ minimizer, every prime $N/7<p\le N$ has positive mass. Throughout $(N/6,N/5]$ the two optimizers therefore make opposite support decisions.
\end{enumerate}
\end{theorem}

One can remove all primes in $(N/6,N/5]$ simultaneously without changing the $K^T$ minimum. The prime number theorem gives $(1/30+o(1))\pi(N)$ such primes. Removing the same set strictly increases the $K^V$ minimum for all sufficiently large $N$, by uniqueness. All assertions about zero coordinates here are exact.

\subsection{Arithmetic decomposition of the stationarity equations}
The matrix $K^T_N$ is positive definite. Indeed,
\[
 z^{\mathsf T}K^T_Nz=\sum_{d\le N}\varphi(d)
    \left(\sum_{d\mid m,\ m\le N}\frac{z_m}{\sqrt m}\right)^2 .
\]
Vanishing forces every $z_m=0$, by descending through the indices $d$. For $K^V_N$ use instead
\[
 z^{\mathsf T}K^V_Nz=\sum_{d\le N}\varphi(d)\int_0^\infty
    \left(\sum_{d\mid m,\ m\le N}z_m e^{-mt}\right)^2\,dt.
\]
Already the $d=1$ term is positive for a nonzero vector by linear independence of distinct powers of $e^{-t}$. The minimizing probability vector $w^{(N)}$ is consequently unique.

The first-order condition at $w=w^{(N)}$ is $(K_Nw)_i\ge\lambda_N$, with equality if $w_i>0$. This follows by moving probability toward any single coordinate and then averaging these inequalities with $w$. Equivalently $v^{(N)}=w^{(N)}/\lambda_N$ is the unique solution of
\begin{equation}\label{eq:gal-obstacle}
 v\ge0,\qquad K_Nv\ge\boldsymbol1,\qquad v_i((K_Nv)_i-1)=0 .
\end{equation}
It maximizes $2\boldsymbol1^{\mathsf T}v-v^{\mathsf T}K_Nv$ on the nonnegative orthant, with maximum $C_N=\boldsymbol1^{\mathsf T}v^{(N)}$. These rescaled coordinates can be optimized on one set while holding the complement fixed.

\begin{theorem}[Multiples of a large prime]\label{thm:gal-block}
Fix $k\ge1$, and take primes $p$ with $\lfloor N/p\rfloor=k$. Uniformly over those primes,
\[
 \left(v^{(N)}_{p},v^{(N)}_{2p},\dots,v^{(N)}_{kp}\right)
 =v^{(k)}+
 \begin{cases}
 O_k((N\lambda_N)^{-1/2}),&K=K^T,\\
 O_k((N\lambda_N)^{-1}),&K=K^V .
 \end{cases}
\]
The error denotes a bound in Euclidean norm. A positive coordinate of $v^{(k)}$ remains positive in the corresponding block for all sufficiently large $N$. If $(K_kv^{(k)})_j>1$, the corresponding coordinate $v^{(N)}_{jp}$ equals zero for all sufficiently large $N$.
\end{theorem}
\begin{proof}
Both kernels obey $K(pi,pj)=K(i,j)$. Put $b_j=\sum_{m\le N,\,p\nmid m}K(pj,m)v^{(N)}_m$ for $1\le j\le k$. Since $p\nmid m$ gives $\gcd(pj,m)=\gcd(j,m)\le j$, the $K^T$ estimate is
\[
 0\le b_j\le \sqrt{j/p}\,\frac{1}{\lambda_N}
       \sum_{m\le N}\frac{w^{(N)}_m}{\sqrt m}
 \le\sqrt{\frac{j}{p\lambda_N}}.
\]
The last inequality uses $\lambda_N\ge(\sum w^{(N)}_m/\sqrt m)^2$. For $K^V$ we have $K^V(pj,m)\le 1/p$, hence $b_j\le1/(p\lambda_N)$. As $p>N/(k+1)$, these bounds tend to zero uniformly by \cite[Theorem~1.1]{dmtgal}.

With the complement fixed, $u=(v^{(N)}_{jp})_{j\le k}$ minimizes $u^{\mathsf T}K_ku/2-(\boldsymbol1-b)^{\mathsf T}u$ on $u\ge0$. Let $u_0=v^{(k)}$ be the solution for $b=0$. Test the two variational inequalities at each other's solution and add. This gives
\[
 \lambda_{\min}(K_k)\|u-u_0\|_2^2
 \le (u-u_0)^{\mathsf T}K_k(u-u_0)
 \le b^{\mathsf T}(u_0-u)\le\|b\|_2\|u-u_0\|_2.
\]
This proves the error estimate. A positive limit coordinate stays positive. If $(K_ku_0)_j>1$, then $(K_ku+b)_j>1$ eventually, forcing $u_j=0$ by \eqref{eq:gal-obstacle}.
\end{proof}

\subsection{Six finite systems and their strict inequalities}
The table gives the supports $S$ in \eqref{eq:gal-obstacle}. The columns labelled $m$ give certified lower bounds $m/1000$ for the smallest positive coordinate after the rescaling $t_j=v_j/\sqrt j$ in the $T$ system, or in the original coordinates $t_j=v_j$ in the $V$ system.
\[
\begin{array}{c|l r|l r}
k&S_T&m_T&S_V&m_V\\ \hline
1&1&1000&1&2000\\
2&1,2&414&1,2&1200\\
3&1,2,3&219&1,2,3&679\\
4&1,2,3,4&121&1,2,3,4&177\\
5&2,3,4,5&88&1,2,3,4,5&112\\
6&1,3,4,5,6&39&1,3,4,5,6&476
\end{array}
\]
In the $T$ system the active equations are $G_{SS}t_S=(\sqrt i)_{i\in S}$, where $G_{ij}=\gcd(i,j)$. For $k=5$ the inactive slack $\sum_j G_{1j}t_j-1$ is
\[
 -1+(4\sqrt2+2\sqrt3+\sqrt5)/11>32/1000 .
\]
For $k=6$ the slack at 2 in the rescaled equations is
\[
 \sum_jG_{2j}t_j-\sqrt2=(5-4\sqrt2-\sqrt3+\sqrt6)/4>15/1000 .
\]
For the $V$ system solve $K^V_{SS}t_S=\boldsymbol1$ over the rationals. Its only inactive slack in this table is at 2 for $k=6$, and exceeds $47/1000$. The following exact calculation verifies the lower bounds and all inactive slacks. The square-root enclosures follow by squaring integers.
\begin{verbatim}
from fractions import Fraction as F
from math import gcd, isqrt

def inverse(A):
    n = len(A)
    L = [row[:] + [F(i == j) for j in range(n)]
         for i, row in enumerate(A)]
    for j in range(n):             # positive definite principal matrices
        L[j] = [t/L[j][j] for t in L[j]]
        for i in range(n):
            if i != j:
                c = L[i][j]
                L[i] = [a-c*b for a,b in zip(L[i],L[j])]
    return [row[n:] for row in L]

def root(i):
    Q = 10**6
    a = isqrt(i*Q*Q)
    return F(a,Q), F(a + (a*a != i*Q*Q),Q)

def lower(row, intervals):
    return sum(c*(lo if c >= 0 else hi)
               for c,(lo,hi) in zip(row,intervals))

for T in (True,False):
    bounds = ([1000,414,219,121,88,39] if T else
              [2000,1200,679,177,112,476])
    for k in range(1,7):
        I = list(range(1,k+1))
        S = [i for i in I if not ((T and k==5 and i==1)
                                   or (k==6 and i==2))]
        A = [[F(gcd(i,j),1 if T else i+j) for j in S] for i in I]
        P = inverse([A[i-1] for i in S])
        r = [root(i) if T else (F(1),F(1)) for i in I]
        rhs = [r[j-1] for j in S]
        assert all(lower(row,rhs) >= F(bounds[k-1],1000)
                   for row in P)
        for i in set(I)-set(S):
            row = [sum(a*b for a,b in zip(A[i-1],column))
                   for column in zip(*P)]
            bound = 32 if k==5 else (15 if T else 47)
            assert lower(row,rhs)-r[i-1][1] > F(bound,1000)
\end{verbatim}
Extending each solution by zero gives \eqref{eq:gal-obstacle}, so positive definiteness proves it is the optimizer. The interval assertions in Theorem~\ref{thm:gal-primes} now follow from Theorem~\ref{thm:gal-block}. The proof gives an asymptotic assertion uniform in $p$ and does not specify a numerical starting value for $N$.

\subsection{Single-integer increments of the optimum}
\begin{theorem}\label{thm:gal-jumps}
For either kernel, every fixed $k\ge1$ satisfies
\[
 \lim_{\substack{p\to\infty\\p\ {\rm prime}}}
   \bigl(C_{kp}-C_{kp-1}\bigr)=C_k-C_{k-1}.
\]
In particular the successive differences $C_N-C_{N-1}$ do not converge. Their limits along primes and along twice a prime are $1$ and $3-2\sqrt2$ for $K^T$, and $2$ and $2/5$ for $K^V$.
\end{theorem}
\begin{proof}
Let $B$ be the set of integers up to $kp$ outside the multiples of $p$. It is the same for $kp-1$. For any fixed outside vector $z\ge0$ on $B$, set $b_j=\sum_{m\in B}K(pj,m)z_m$, $1\le j\le k$. After optimizing the first $r$ multiples of $p$, their contribution to the orthant maximum is
\[
 J_r(b)=\max_{u\ge0}\{2(\boldsymbol1-b_{1:r})^{\mathsf T}u-u^{\mathsf T}K_ru\},
 \qquad J_0=0 .
\]
The value of the full objective is the outside-only objective plus $J_r(b)$. We have $J_r(0)=C_r$ and $J_r(b)\to C_r$ when $b\ge0$ tends to zero. Inserting the unperturbed optimizer gives $C_r-2b_{1:r}^{\mathsf T}v^{(r)}\le J_r(b)\le C_r$. For $r=0$ use $C_0=0$.

Take the outside part of the full optimizer for $kp-1$, giving $b^-$. Keeping this outside vector as a trial vector for $kp$ gives $C_{kp}-C_{kp-1}\ge J_k(b^-)-J_{k-1}(b^-)$. Using the outside part of the optimizer for $kp$, giving $b^+$, yields the reverse bound with $b^+$. The off-block estimates in the preceding proof apply to either outside vector for all $j\le k$, including $j=k$ even at size $kp-1$. Both $b^-$ and $b^+$ tend to zero. The stated limit follows. Finally $C_1^T=1$, $C_2^T=2/(1+1/\sqrt2)$, $C_1^V=2$, and $C_2^V=12/5$, by the one- and two-variable systems.
\end{proof}


\begin{corollary}\label{cor:gal-jump-limits}
Put $d_n=C_n-C_{n-1}$. For $K^T$ one has $\liminf d_n=0$ and $\limsup d_n=1$. For $K^V$ these values are $0$ and $2$. For either kernel, $\{n:d_n\ge\epsilon\}$ has natural density zero for every $\epsilon>0$. The set of subsequential limits of $(d_n)$ is the closure of $\{d_k:k\ge1\}$.
\end{corollary}
\begin{proof}
Setting the new coordinate to zero proves $d_n\ge0$. With the first $n-1$ coordinates held fixed, the gain from an additional coordinate $t\ge0$ is at most $2t-K(n,n)t^2$, since all mixed terms in the quadratic form are nonnegative. Maximizing gives $d_n\le1/K(n,n)$. Theorem~\ref{thm:gal-jumps} with $k=1$ attains this upper bound in the limit. Also
\[
 \frac1X\#\{n\le X:d_n\ge\epsilon\}
 \le\frac{C_X}{\epsilon X}=\frac1{\epsilon X\lambda_X}\longrightarrow0
 \qquad(X\to\infty,\ X\in\mathbb N).
\]
This implies the asserted lower limit. Each $d_k$ is a subsequential limit by the same theorem. The cluster set is closed and contained in the closure of the range, giving equality.
\end{proof}

The support intervals and the increment law describe the exact optimization problems in \cite{dmtgal}. They leave the secondary factors in the growth estimates for $T_N,V_N$ undetermined.

\section{The original distance energy on the integers}\label{sec:metric}
Join $m$ to $mp$ whenever $p$ is prime and both numbers belong to $[N]=\{1,\ldots,N\}$. A path through $\gcd(a,b)$ shows that the graph distance is
\[
 d(a,b)=\sum_p|v_p(a)-v_p(b)|=\Omega(ab/\gcd(a,b)^2).
\]
Each step changes just one prime exponent by one, so the path has minimum length. This prime-factor distance is also described in \cite{dist}. For independent $X,Y$ with common law $w$ on $[N]$, define the exact objective from the graph manuscript by
\[
 F_N=\max_w\left\{\E d(X,Y)+2\sum_{m\le N}w_m^2\right\}.
\]
\begin{theorem}[Six primes determine the leading energy]
Set $P=\{2,3,5,7,11,13\}$, $\lambda=4/\log 30030$, and $s_p=(1-\lambda\log p)/2$ for $p\in P$, with $s_p=0$ elsewhere. Then
\[
 F_N=C\log N+O(1),\qquad C=\sum_{p\in P}\frac{2s_p(1-s_p)}{\log p}=1.3715863063\ldots .
\]
For any maximizing law (also for any law within a fixed additive constant of the maximum),
\[
 \sum_{p>13}(\log p)\,\E v_p(X)=O(1),\qquad
 \frac{\log p}{\log N}\E v_p(X)=s_p+O((\log N)^{-1/2})\quad (p\in P).
\]
\end{theorem}
\begin{proof}
Let $L=\log N$ and $u_p=(\log p)\E v_p(X)/L$. The $u_p$ are nonnegative and their sum is at most one. Put $t_{p,j}=\Pr(p^j\mid X)$, $1\le j\le k_p=\lfloor L/\log p\rfloor$ for $p\le N$. Then $\sum_j t_{p,j}=\E v_p(X)$ and
\[
 \E d(X,Y)=2\sum_{p,j}t_{p,j}(1-t_{p,j})
 \le L\Phi(u),\qquad
 \Phi(u)=\sum_p\frac{2u_p(1-u_p)}{\log p}.
\]
The inequality follows from $\sum_j t_{p,j}^2\ge(\sum_jt_{p,j})^2/k_p$ and $k_p\le L/\log p$. The collision term is at most two.
 
Since $13^4<30030<17^4$, the positive coordinates $s_p$ are exactly those in $P$. They sum to one. On those coordinates the derivative of $\Phi$ at $s$ is $2\lambda$. On a zero coordinate it is $2/\log p<2\lambda$. More precisely,
\begin{align*}
 C-\Phi(u)={}&2\lambda\left(1-\sum_p u_p\right)
 +\sum_{p>13}\left(2\lambda-\frac{2}{\log p}\right)u_p
 +\sum_p\frac{2(u_p-s_p)^2}{\log p}.
\end{align*}
Thus the relaxed optimizer and its support are explicit. For $N\ge13$ put mass $s_p$ on $p^{k_p}$ for $p\in P$. For each active prime, $t_{p,j}=s_p$ at every positive level, so the expected distance is $2\sum_{p\in P}k_p s_p(1-s_p)=CL+O(1)$. This proves the estimate for the discrete graph. A near-optimal law has $L(C-\Phi(u))=O(1)$. Above $13$ the coefficients in the middle sum are at least $2\lambda-2/\log17>0$, and the last sum controls the six stated errors.
\end{proof}
The six numbers $s_p$ describe shares of the expected logarithmic size. The theorem permits finite-$N$ optimizers to give some mass to mixed integers and to primes above $13$. This example keeps the graph-distance kernel exactly. Its main estimate uses stationarity after passing to prime-exponent marginals.

\section{Equal pairwise least common multiples on two prime-factor layers}
\label{sec:lcm-transfer}
Erd\H{o}s asked for the largest reciprocal sum of a subset of $[1,N]$ containing
no $k$ distinct integers whose pairwise least common multiples are all equal
\cite{erdos-lcm}. Tang and Zhang connect the unrestricted problem with
sunflower-free capacity \cite{tangzhang-lcm}. We compute the semiprime
restriction with absolute error tending to zero and identify the coefficient
on the next layer with the ordinary Tur\'an density of the tetrahedron.

Write $\Omega(m)$ for the number of prime factors with multiplicity, and set
\[
 S=S(N)=\sum_{p\le N}\frac1p,\qquad P(2)=\sum_p\frac1{p^2}.
\]
For $j=2,3$ let $L_j(N)$ maximize $\sum_{a\in A}1/a$ over sets
$A\subseteq\{m\le N:\Omega(m)=j\}$ containing no $j+1$ distinct
integers with the same pairwise least common multiple.

\paragraph{The forbidden configurations.}
If $a_1,\ldots,a_{j+1}$ form such a configuration and $L$ is their
common pairwise least common multiple, the integers $L/a_i$ are pairwise
coprime. Each has $t=\Omega(L)-j$ prime factors, where $t\ge1$ since
the $a_i$ are distinct. Their product divides $L$, so
$(j+1)t\le j+t$. Thus $t=1$ and the $j+1$ deficits are distinct primes.
Their product exhausts the prime factors of $L$. In particular $L$ is
squarefree and the $a_i$ correspond exactly to all $j$-subsets of a
$(j+1)$-element set of primes. The converse follows by taking pairwise
unions. All nonsquarefree members of the layer may therefore be included.

\begin{theorem}[The semiprime restriction]\label{thm:lcm-semiprimes}
As $N\to\infty$,
\[
 L_2(N)=\frac{S(N)^2}{4}+P(2)+o(1).
\]
\end{theorem}
\begin{proof}
Represent each selected product $pq$, $p<q$, by an edge between primes.
The graph is triangle-free and has edge mass
$E=\sum_{pq\text{ edge}}x_px_q$ for $x_p=1/p$. Holding the graph
fixed, temporarily vary these masses on $\sum x_p=S$. If positive masses
at $u,v$ are nonadjacent, move the entire mass of the one with smaller
weighted neighbor sum onto the other. The edge mass is linear in the
amount moved because $uv$ is absent. It cannot decrease. Each operation
deletes one positive coordinate, so the final support is a clique of size
at most two. Hence $E\le S^2/4$. This is the quadratic simplex argument
of Motzkin and Straus \cite{motzkin-lcm}; the moved masses certify an
upper bound, while the arithmetic masses in the construction remain $1/p$.
For a precise link with stationarity--deletion, normalize the total mass to
one and let $E_*$ be the edge mass at a simplex maximizer. The neighbor
sums on its support all equal $2E_*$ by stationarity and double counting.
Deleting a coordinate of mass $t<1$ and normalizing the others therefore
gives the exact identity $(1-t)^2E_{-}=(1-2t)E_*$. Here $E_{-}$ is
the edge mass at the renormalized vector, with the deleted vertex absent.

For the lower bound let $y$ be the first prime with $S(y)\ge S/2$.
Then $0\le S(y)-S/2\le1/y$. The quantitative prime-reciprocal
estimate of Mertens,
$S(x)=\log\log x+B+O(1/\log x)$, gives
$\log y\asymp\sqrt{\log N}$. Take all cross products $pq\le N$ with
$p\le y<q\le N$. Before the product cutoff their mass equals
$S^2/4-(S(y)-S/2)^2$. The cutoff removes at most
\[
 S(y)\,[S(N)-S(N/y)]
   \ll (\log\log N)\frac{1+\log y}{\log N}=o(1),
\]
by the same Mertens estimate. The bipartite graph is triangle-free.
Adding every $p^2\le N$ contributes $P(2)+o(1)$, which also bounds
the square contribution from above.
\end{proof}

The deficit has a structural interpretation. For a triangle-free selection
put $\delta=S^2/4-E$, take a vertex of maximum weighted neighbor sum $a$,
and let $U$ be its neighborhood. Write $I$ for the edge mass internal
to the complement of $U$, and $J$ for the missing cross-edge mass relative
to the complete bipartite graph with parts $U,U^c$ on all primes up to $N$.
Since $U$ is independent, summing weighted degrees over $U^c$ gives
$a(S-a)-J+2I\le a(S-a)$. Consequently
\[
  (a-S/2)^2+(I+J)/3\le\delta.
\]
Thus any sequence of selections within $o(1)$ of the optimal reciprocal sum
on this layer differs from such bipartite prime products by $o(1)$
reciprocal mass. The comparison includes oversized cross products as
missing. This stability estimate concerns squarefree products.

\begin{theorem}[The three-prime-factor restriction]\label{thm:lcm-tetra}
Let $\pi_4$ be the ordinary Tur\'an density of the complete
three-uniform hypergraph on four vertices. Then
\[
 L_3(N)=\frac{\pi_4}{6} S(N)^3+O\bigl(S(N)\bigr).
\]
In particular $\lim_{N\to\infty}L_3(N)/S(N)^3=5/54$ is equivalent to
Tur\'an's tetrahedron conjecture $\pi_4=5/9$ \cite{turan-lcm}.
\end{theorem}
\begin{proof}
Represent squarefree products by triples of primes. The forbidden
configuration is the tetrahedron. Let $Q$ denote the total mass $1/(pqr)$
of all products $p<q<r$, $pqr\le N$, and let $F$ be the maximum admissible
mass on these triples. We claim
\[
 \pi_4 Q\le F\le \pi_4 S^3/6.
\]
For the upper bound, replace prime $p$ in an optimal hypergraph by a class
of proportion $1/(pS)$, and replace each selected triple by all triples
meeting its three classes. These proportions are rational, so arbitrarily
large exact integer class sizes exist. The resulting hypergraph remains
tetrahedron-free: two vertices from one class occur together in no edge,
whereas every pair of tetrahedron vertices must occur in an edge.
Taking the limit of its edge density yields $6F/S^3\le\pi_4$.
Conversely take an extremal tetrahedron-free hypergraph on $m$ vertices.
Label each prime independently and uniformly by its vertices, and keep
legal prime triples whose labels form an edge. The same pair argument
makes the selection tetrahedron-free. Its expected mass is
$6\operatorname{ex}(m,K_4^{(3)})Q/m^3$. Letting $m\to\infty$ gives the
lower bound by the definition of Tur\'an density.

It remains to estimate $Q$. Repeated primes in the expansion of $S^3$
have mass $O(P(2)S)$. For distinct triples exceeding the cutoff, write
$p\ge q\ge r$; necessarily $pq^2>N$. Such ordered pairs $(p,q)$ have
total mass $O(1)$. Indeed for $q>N^{1/3}$ both factors lie in a prime
interval with bounded reciprocal sum. For $q\le N^{1/3}$, Mertens'
estimate bounds the sum over $N/q^2<p\le N$ by
$O((1+\log q)/\log N)$. Summing with weights $1/q$ gives $O(1)$
by $\sum_{q\le x}(\log q)/q=O(\log x)$ over primes. The third factor
costs at most $S$, so $Q=S^3/6+O(S)$. Finally all allowed nonsquarefree
numbers on this layer have total mass $O(P(2)S)$, by writing them
as $p^2q$. This proves the estimate.
\end{proof}

These layer computations leave the unrestricted growth problem open.
In particular a power of $\log\log N$ on a fixed layer does not improve
the power-of-$\log N$ lower bounds for the unrestricted problem.

\subsection{The cost of the product cutoff and the counting version}
The host graph of semiprimes has an extra ordering property. Write
$X=\sqrt N$. The primes up to $X$ support all pairs, primes above $X$
support none among themselves, and the possible small neighbors of $q>X$
are precisely the initial prime segment $p\le N/q$. This determines the
extremal counting problem at every $N$ and the first correction term in
Theorem~\ref{thm:lcm-semiprimes}.

\begin{lemma}\label{lem:lcm-prefix}
Give each prime $p\le N$ a positive mass $w_p$. Put
$W(z)=\sum_{p\le z}w_p$, for $z\le N$, and $W_0=W(X)$.
The selected squarefree edge mass of any admissible semiprime family
is bounded by the maximum on $0\le a\le W_0$ of
\[
 \Phi(a)=a(W_0-a)+\sum_{X<q\le N}w_q\min\{a,W(N/q)\}.
\]
For every $a=W(y)$, where $y$ is a prime up to $X$, or $a=0$,
the value $\Phi(a)$ is attained by taking all legal edges across the cut
between the initial prime segment and its complement.
\end{lemma}
\begin{proof}
Let $a$ be the largest mass of an independent set in the selected graph on
small primes. All its weighted degrees are at most $a$, since the selected
graph is triangle-free. Summing degrees with vertex weights over the
complement of an independent set of mass $a$ counts every selected core edge
at least once. Hence its edge mass is at most $a(W_0-a)$. The small
neighbors of a selected large prime form an independent set within its
allowed prefix, with mass at most $\min\{a,W(N/q)\}$. All inequalities in
the stated formula for a prefix cut follow by counting crossing edges
directly.
\end{proof}

\begin{theorem}[First correction term on semiprimes]\label{thm:lcm-correction}
For the prime reciprocal constant $B$ defined by
$S(x)=\log\log x+B+o(1)$, one has
\[
 L_2(N)=\frac{S(N)^2}{4}+P(2)
       -\frac{e^{-B/2}+o(1)}{\sqrt{\log N}}.
\]
\end{theorem}
\begin{proof}
Use $w_p=1/p$ in Lemma~\ref{lem:lcm-prefix}, and set $T=\log N$.
For $a$ in its range write
$\Phi(a)=a(S-a)-J(a)$, where
$J(a)=\sum_{q>X}q^{-1}(a-S(N/q))_+\ge0$ and the sum is over $q\le N$.
If $a=S/2+o(1)$, take the prime $y$ where this prefix mass ends:
$a=S(y)- (1-\theta)/y$ with $0\le\theta\le1$.
Mertens' formula gives $\log y=(e^{-B/2}+o(1))\sqrt T$.
Exchanging the sums over incompatible cross pairs yields
\[
 J(a)=\sum_{p<y}\frac{S(N)-S(N/p)}p+
        \frac{\theta}{y}[S(N)-S(N/y)].
\]
Uniformly for $p\le y$, the quantitative Mertens estimate gives
$S(N)-S(N/p)=(\log p)/T+O((\log p)^2/T^2+1/T)$.
Using $\sum_{p\le y}(\log p)/p=\log y+O(1)$, the quadratic
remainder after summation is $O((\log y)^2/T^2)$, and the constant
remainder is $O(S(y)/T)$. The fractional endpoint changes the sum by
at most $S/y$. Consequently
\[
 J(a)=\frac{e^{-B/2}+o(1)}{\sqrt T}
 \qquad\text{whenever } a=S/2+o(1).
\]

A maximizing $a$ for $\Phi$ has $a=S/2+o(1)$: Theorem~\ref{thm:lcm-semiprimes}
gives an attainable squarefree mass $S^2/4-o(1)$, whereas
$\Phi(a)\le S^2/4-(a-S/2)^2$. The formula for $J$ bounds the maximum
from above. For the lower bound take the first actual prime prefix reaching
$S/2$; its squared mass imbalance is at most $1/y^2=o(T^{-1/2})$,
and its value of $\Phi$ is attained. Finally the tail in the prime-square
contribution is $O(N^{-1/2})$.
\end{proof}

\begin{theorem}[The semiprime restriction of Erd\H{o}s's counting problem]
\label{thm:lcm-count}
For every positive integer $N$, among all semiprimes up to $N$ the maximum number selectable without
three equal pairwise least common multiples is attained by all prime
squares together with the legal squarefree products across a single
prime-size threshold. With $m=\pi(\sqrt N)$ the exact maximum is
\[
 m+\max_{0\le h\le m,\ h\in\mathbb Z}
  \left[h(m-h)+\sum_{\sqrt N<q\le N}\min\{h,\pi(N/q)\}\right].
\]
As $N\to\infty$, the minimum number of semiprimes that must be omitted
is $(\tfrac12\log 2+o(1))m^2$, or equivalently
$(2\log 2+o(1))N/(\log N)^2$.
\end{theorem}
\begin{proof}
Apply Lemma~\ref{lem:lcm-prefix} with all masses one. The actual
independence number is an integer, and every integer prefix is available;
all $m$ squares may be added. If $p_j$ is the $j$th prime, the forward
difference of the expression in square brackets at $h<m$ equals
$\pi(N/p_{h+1})-2h-1$ and is strictly decreasing. For
$h/m\to c\in(0,1)$ the prime number theorem gives
$p_{h+1}/\sqrt N\to c$ and the forward difference divided by $m$ tends
to $c^{-1}-2c$. Thus every maximizing $h$ has $h/m\to t=1/\sqrt2$
and $y=p_h$ satisfies $y/\sqrt N\to t$.

The excluded squarefree products have both primes at most $y$ or both
above it. Their number equals
\[
 \binom h2+\sum_{y<p\le\sqrt N}\bigl(\pi(N/p)-\pi(p)\bigr).
\]
The prime number theorem gives $\pi(u\sqrt N)/m\to u$ uniformly on each
compact subinterval of $(0,\infty)$ (by monotonicity and continuity of
the limit). Thus the last sum divided by $m^2$ tends, by Riemann
summation over scaled primes, to $\int_t^1(u^{-1}-u)\,du$. Including
the first term gives $t^2-\tfrac12-\log t=\tfrac12\log 2$.
\end{proof}

An independent exact search branching on the edges of each remaining triangle verified the counting formula for every $1\le N\le120$.

The counting question for unrestricted integers is Erd\H{o}s problem 536,
while the reciprocal-sum question is problem 856 \cite{erdos-lcm}. Their respective
semiprime restrictions have different extremizing scales:
$\log y\sim e^{-B/2}\sqrt{\log N}$ in the harmonic construction and
$y\sim\sqrt{N/2}$ for an optimal cardinality threshold.

\section{The constant in Erd\H{o}s's equal-gcd problem}\label{sec:gcd-sunflower}
\begingroup\setlength{\emergencystretch}{1.5em}
Fix an integer $r\ge3$. Let $f_r(N)$ be the largest size of a subset of
$[1,N]\cap\mathbb Z$ containing no $r$ distinct integers whose pairwise
greatest common divisors are all equal. Let $F_r(k)$ be the largest size
of a family of distinct $k$-element sets containing no $r$ distinct
members whose pairwise intersections are all equal (an $r$-sunflower).
The ground set is unrestricted. The sunflower lemma gives
$F_r(k)\le k!(r-1)^k$. Put
\[
 C_r=\sup_{k\ge1} F_r(k)^{1/k}\in[r-1,\infty],\qquad
 B(N)=\frac{\log N}{\log\log N}.
\]
The lower endpoint follows from $F_r(1)=r-1$. All logarithms are natural.

\begin{theorem}\label{thm:gcd-sunflower-constant}
For each fixed $r\ge3$, the following equality holds in the extended real
numbers:
\[
 \lim_{N\to\infty}\frac{\log f_r(N)}{B(N)}=\log C_r.
\]
In particular, $f_r(N)\le\exp(O_r(B(N)))$ holds if and only if the
Erd\H{o}s--Rado sunflower conjecture holds for this $r$. If $C_r=\infty$,
the displayed quotient tends to $+\infty$.
\end{theorem}

\begin{proof}
First suppose $C_r<\infty$. Write $L=\log N$, $\ell=\log\log N$ and
$y=L/\ell^3$. Split $n\le N$ uniquely as $ab$, assigning all prime powers
with prime base $p\le y$ to $a$ and the others to $b$. Since every prime
exponent is at most $L/\log 2$, the number of possible values of $a$ is at
most $(1+L/\log 2)^{\lfloor y\rfloor}=\exp(O(L/\ell^2))$. The number
$j=\Omega(b)$ of prime factors counted with multiplicity is at most
$K=\lfloor L/\log y\rfloor=(1+o(1))B(N)$. Encode $b$ by the $j$-set
\[
 S(b)=\{(p,t):1\le t\le v_p(b)\}.
\]
Then $S(b_1)\cap S(b_2)=S(\gcd(b_1,b_2))$. At fixed $a$ and $j$,
the encoding is injective, and an $r$-sunflower would give $r$ original
integers with equal pairwise gcds, since $\gcd(ab_1,ab_2)=a\gcd(b_1,b_2)$.
Thus an admissible set has at most $F_r(j)\le C_r^j$ integers in each
such class, with the bound $1$ for $j=0$. Consequently
\[
 f_r(N)\le\exp(O(L/\ell^2))(K+1)C_r^K.
\]
Dividing logarithms by $B(N)$ gives the required upper bound.

For the lower bound use Abbott's finite-block construction
\cite[p.~159, (7)]{gcdAbbott66}. Fix $k$ and a sunflower-free $k$-uniform
family $\mathcal F$ of size $F_r(k)$ on $h$ elements. On $m$ disjoint
copies of its ground set take all unions obtained by choosing one member
of $\mathcal F$ from each copy. If $r$ such unions form a sunflower,
then in each copy the coordinate choices have constant pairwise
intersection. They cannot all be distinct. If two are equal to $A$,
all others contain $A$ and have size $k$, so that every choice in this
copy equals $A$. Doing this in every copy makes the $r$ unions identical.
The resulting $F_r(k)^m$ distinct unions are therefore sunflower-free.

Represent the $s=hm$ elements by the first $s$ primes, and each union by
the product of its primes. We obtain admissible squarefree integers with
exactly $km$ prime factors, bounded by $p_s^{km}$, where $p_s$ is the
$s$th prime. Fix $0<\epsilon<1$ and take
$m=\lfloor(1-\epsilon)B(N)/k\rfloor$. The estimate $p_s=O(s\log(2s))$
gives $p_s^{km}\le N^{1-\epsilon/2}$ for all sufficiently large $N$.
Hence the lower limit is at least $(1-\epsilon)\log F_r(k)/k$ for every
fixed $k$. Send $\epsilon$ to zero and take the supremum over $k$. This
also proves divergence to $+\infty$ when that supremum is infinite.
\end{proof}

Erd\H{o}s explicitly asked about this normalized limit in 1964 and
sketched the sunflower-to-arithmetic implication
\cite[pp.~645--646, (7)--(8)]{gcdErdos64}. His forcing threshold exceeds
our extremal size by one. Abbott's bound already gives the reverse
implication, including divergence when $C_r=\infty$. In 1973 Erd\H{o}s
reported a concern about applying the ordinary sunflower conjecture to
repeated prime factors \cite[Sec.~4]{gcdErdos73}. The prime-exponent
encoding above handles those factors and the choice of $y$ makes the
smooth-part cost $o(B(N))$ in the logarithm. Thus the formula identifies
the coefficient in terms of the set problem. Its finiteness remains the
open fixed-$r$ sunflower conjecture \cite{gcdRao26}.

\paragraph{The additional coprimality condition in Erd\H{o}s's 1973 formulation.}
Call a tuple \emph{proper} if its pairwise gcds all equal $d$ and
$\gcd(n_i/d,d)=1$ for each member. Thus every prime dividing two members
must divide all members to exactly the same exponent. Write $G_r(k)$ for
the largest size of a set of distinct positive integers with
$\Omega(n)=k$ containing no proper $r$-tuple. Erd\H{o}s's conjecture
(4.3) asks for $G_r(k)\le D_r^k$ \cite[p.~123]{gcdErdos73}.

\begin{theorem}\label{thm:proper-gcd-capacity}
For every fixed $r\ge3$ one has
\[
 \lim_{k\to\infty}G_r(k)^{1/k}=1+C_r
\]
in the extended real numbers. In particular, conjectures (4.2) and (4.3)
of Erd\H{o}s \cite{gcdErdos73} are equivalent for each $r$.
\end{theorem}
\begin{proof}
For a lower bound, realize an extremal $j$-uniform family as prime
supports and give each of its supports all possible positive exponents
summing to $k$. In a proper tuple the supports have constant pairwise
intersection. If two supports agree, the others must contain that support
and hence equal it, since all have size $j$. Properness then forces all
integers to coincide. If all supports are distinct they would form a
sunflower. Thus for $1\le j\le k$,
\[
 G_r(k)\ge F_r(j)\binom{k-1}{j-1}.
\]
Abbott's product construction gives $F_r(i+j)\ge F_r(i)F_r(j)$
\cite[Thm.~1]{gcdAbbott66}, so Fekete's lemma gives
$F_r(j)^{1/j}\to C_r$, allowing infinity. For $1<c<C_r$ we then have
$F_r(j)\ge c^j$ for all sufficiently large $j$. The full sum of
$c^j\binom{k-1}{j-1}$ is $c(1+c)^{k-1}$. Discarding finitely many
fixed $j$ costs only a polynomial in $k$, and at least one of the
remaining terms is at least their sum divided by $k$. Taking $k$th roots
and then the supremum over $c$ gives the required lower limit.

If $C_r<\infty$, independently assign each relevant prime a label
$a\ge1$ with probability $C_r/(1+C_r)^a$. For a finite admissible set
in the slice $\Omega(n)=k$, $\omega(n)=j$, retain $n$ if every prime
in its support received its exponent in $n$. The survival probability is
$C_r^j/(1+C_r)^k$. On the retained integers the map to their supports is
injective, and a sunflower of supports gives a proper tuple since a
shared prime carries a single label. The expected number retained is
therefore at most $F_r(j)\le C_r^j$. The original slice has size at most
$(1+C_r)^k$, and summing over $1\le j\le k$ bounds $G_r(k)$ by
$k(1+C_r)^k$. When $C_r=\infty$ the lower bound suffices. Finiteness
of $G_r(k)$ for each individual $k$ follows by the same argument with
label probabilities $2^{-a}$, giving $G_r(k)\le2^k\sum_{j\le k}F_r(j)$,
where the sum runs from $1$ to $k$.
\end{proof}

Write $f_r^{\mathrm{prop}}(N)$ for the corresponding extremum in $[1,N]$. Then
\[
 \lim_{N\to\infty}\frac{\log f_r^{\mathrm{prop}}(N)}{B(N)}=\log(1+C_r).
\]
For finite $C_r$, use the smooth-factor decomposition from
Theorem~\ref{thm:gcd-sunflower-constant}. A proper tuple of rough factors
remains proper after multiplication by the common smooth factor, because
their prime supports are disjoint from that factor. Each rough slice of
rank $j\ge1$ therefore contains at most $G_r(j)\le j(1+C_r)^j$ elements.
For the lower bound take an extremal $\Omega(n)=k$ family and form the
products across $m$ copies on disjoint primes, as in that theorem. In a
proper tuple the entries in each copy satisfy the proper equalities.
If two entries coincide, every other entry is divisible by this same
integer and has the same $\Omega$, so all coincide. Distinct entries
would violate the choice of the family. The construction thus gives
$G_r(k)^m$ admissible integers bounded by $p_{hm}^{km}$, where $h$ is
the number of primes in the fixed family. Applying the preceding limit
in $k$ gives the claimed lower limit, also when $C_r=\infty$.

\paragraph{Localization.}
Choose $0<\eta_N\le1$ with $\log(1/\eta_N)=o(B(N))$. Write
$g_r(N,\eta_N)$ for the largest size of an admissible set of squarefree
integers in $[N,(1+\eta_N)N]$ whose members all have the same number of
prime factors. This common number may depend on $N$. Then
\[
 \lim_{N\to\infty}\frac{\log g_r(N,\eta_N)}{B(N)}=\log C_r.
\]
For example, $\eta_N=(\log N)^{-A}$ works for every fixed $A>0$.

\begin{proof}
The upper bound follows from $g_r(N,\eta_N)\le f_r(2N)$ and
$B(2N)/B(N)\to1$. For the lower bound take the squarefree products from
Theorem~\ref{thm:gcd-sunflower-constant}, with $k$ and $\epsilon$ fixed. Partition
$[1,N^{1-\epsilon/2}]$ into geometric intervals $[z,(1+\eta_N/4)z)$.
There are $O(1+L/\eta_N)$ intervals. Some interval thus contains at
least $F_r(k)^m\exp(-o(B(N)))$ products. With $x=N/z$ we have
$N^{\epsilon/2}\le x\le N$. A common multiplier $q$ in
\[
 [x,x+H],\qquad H=x\biggl(\frac{1+\eta_N}{1+\eta_N/4}-1\biggr)
 \asymp x\eta_N,
\]
puts all these products in the required interval. We choose $q$ squarefree
and coprime to $p_1\cdots p_s$, where $s=hm=O(B(N))$.

Indeed, inclusion--exclusion counts at least
$H\prod_{i\le s}(1-1/p_i)-O(2^s)\ge H/(s+1)-O(2^s)$ integers in this
interval coprime to the product, using $p_i\ge i+1$. Among these, those
with a square prime factor have a divisor $p^2$ with $p>p_s$. Their
number is at most
$H\sum_{j>p_s}j^{-2}+O(\sqrt{x})\le H/p_s+O(\sqrt{x})$, since
$x+H\le2x$. As $p_s\ge2s-1$ for $s\ge2$, the number left is
$\gg H/s-O(2^s+\sqrt{x})$, positive for all sufficiently large $N$:
$s\to\infty$, $2^s=N^{o(1)}$, and $H=x^{1-o(1)}$ uniformly for the
permitted $x$. Multiplication by this $q$ preserves squarefreeness and
equal numbers of prime factors. It multiplies every gcd by the same $q$.
Taking logarithms and then letting $\epsilon\downarrow0$ and taking the
supremum over $k$ proves the assertion in both the finite and infinite cases.
\end{proof}
\endgroup

\section*{Which part of the argument transfers?}
\begin{center}\small
\begin{tabular}{@{}>{\raggedright\arraybackslash}p{.25\linewidth}>{\raggedright\arraybackslash}p{.29\linewidth}>{\raggedright\arraybackslash}p{.38\linewidth}@{}}
\toprule
Problem & Quantity that changes & Conclusion and scope\\\midrule
Bhargava selection & One chosen point per unit, with a marginal gain at each nested cluster & One network certifies every maximizing set at every size. Applies to all finite weighted ultra triples, including the $r$-removed version.\\[3pt]
Small Gál sums & Deleting a coordinate of a unique minimizing probability & An entire interval of primes has exact zero mass for one kernel and positive mass for the other.\\[3pt]
Prime-factor distance & Expected exponent at each prime, subject to a logarithmic budget & The six active primes give the sharp leading coefficient of the original distance energy.\\[3pt]
Equal pairwise LCM & Edge mass on pairs of prime vertices & Exact counting optimizer on the semiprime layer. The unrestricted problem allows higher layers.\\[3pt]
Equal pairwise GCD & Intersections of prime-exponent sets & Arithmetic growth constants are expressed through sunflower capacity. Its finiteness is a separate open question.\\
\bottomrule
\end{tabular}
\end{center}
\clearpage
\begin{thebibliography}{99}
\bibitem{gp} D.~Grinberg and F.~Petrov, A greedoid and a matroid inspired by Bhargava's $p$-orderings, \emph{Electron. J. Combin.} \textbf{28} (2021), P3.6. \url{https://doi.org/10.37236/9046}.
\bibitem{gprev} D.~Grinberg and F.~Petrov, extended version, arXiv:1909.01965v5, 13 June 2026. \url{https://arxiv.org/abs/1909.01965v5}.
\bibitem{grin} D.~Grinberg, The Bhargava greedoid as a Gaussian elimination greedoid, \emph{Electron. J. Combin.} \textbf{31} (2024), P2.28. \url{https://doi.org/10.37236/11222}.
\bibitem{kls} B.~Korte, L.~Lovász and R.~Schrader, \emph{Greedoids}, Springer, 1991. \url{https://doi.org/10.1007/978-3-642-58191-5}.
\bibitem{removed} D.~Krachun and R.~Mirgalimova, Strong greedoid structure of $r$-removed $P$-orderings, \emph{Electron. J. Combin.} \textbf{31} (2024), P4.40. \url{https://doi.org/10.37236/12403}.
\bibitem{dmtgal} R.~de la Bretèche, M.~Munsch and G.~Tenenbaum, Small Gál sums and applications, \emph{J. Lond. Math. Soc.} \textbf{103} (2021), 336--352. \url{https://doi.org/10.1112/jlms.12378}.
\bibitem{dist} O.~Leka, Arithmetic distance graphs, transformation kernels, and Liouville spin structures, notes (2026), Sec.~2. \url{https://www.orges-leka.de/arithmetic_distance_graphs_theory.pdf}.
\bibitem{erdos-lcm} T.~F. Bloom (ed.), Erd\H{o}s Problems, nos.~536 and 856. \url{https://www.erdosproblems.com/536} and \url{https://www.erdosproblems.com/856}.
\bibitem{tangzhang-lcm} Q.~Tang and S.~Zhang, Harmonic LCM patterns and sunflower-free capacity, arXiv:2512.20055 (2025). \url{https://arxiv.org/abs/2512.20055}.
\bibitem{motzkin-lcm} T.~S. Motzkin and E.~G. Straus, Maxima for graphs and a new proof of a theorem of Turán, \emph{Canad. J. Math.} \textbf{17} (1965), 533--540.
\bibitem{turan-lcm} P.~Keevash, Hypergraph Turán problems, in \emph{Surveys in Combinatorics 2011}, Cambridge University Press, 2011, 83--140.

\bibitem{gcdErdos64} P.~Erd\H{o}s, On a problem in elementary number theory and a combinatorial problem, \emph{Math. Comp.} \textbf{18} (1964), 644--646. \url{https://doi.org/10.2307/2002950}.
\bibitem{gcdAbbott66} H.~L.~Abbott, Some remarks on a combinatorial theorem of Erd\H{o}s and Rado, \emph{Canad. Math. Bull.} \textbf{9} (1966), 155--160. \url{https://doi.org/10.4153/CMB-1966-019-1}.
\bibitem{gcdErdos73} P.~Erd\H{o}s, Problems and results on combinatorial number theory, in \emph{A Survey of Combinatorial Theory}, North-Holland, 1973, 117--138. \url{https://doi.org/10.1016/B978-0-7204-2262-7.50017-X}.
\bibitem{gcdRao26} A.~Rao, The story of sunflowers, \emph{J. Lond. Math. Soc.} \textbf{113} (2026), e70380. \url{https://doi.org/10.1112/jlms.70380}.

\end{thebibliography}
\end{document}